\documentclass[11pt]{article}
\usepackage[margin=0.82in]{geometry}
\usepackage[T1]{fontenc}
\usepackage{lmodern,microtype,enumitem,tabularx,array,xcolor,hyperref,titlesec}
\definecolor{accent}{HTML}{173B57}
\hypersetup{colorlinks=true,urlcolor=accent,linkcolor=accent}
\setlist{nosep,leftmargin=*}
\setlength{\parindent}{0pt}
\setlength{\parskip}{0.55em}
\titleformat{\section}{\large\bfseries\color{accent}}{}{0pt}{}
\titleformat{\subsection}{\normalsize\bfseries}{}{0pt}{}
\newcommand{\ApplicantName}{Zhiheng Zhang}
\newcommand{\ApplicantEmail}{[CONFIRM EMAIL]}
\newcommand{\ApplicantPhone}{[CONFIRM PHONE]}
\newcommand{\ApplicantGitHub}{\url{https://github.com/zhiheng-zhang-Mera/Quant-Ultra/tree/1988d9a8530da91a8158de864d098ea869098923}}
\pagestyle{plain}
\begin{document}
\begin{center}
{\LARGE\bfseries Research Proposal}\\[0.25em]
{\large \ApplicantName}\\
Drexel University --- Shahin Jabbari
\end{center}
\vspace{0.35em}\hrule\vspace{0.65em}

\textbf{Provisional title:} Auditable Financial Machine Learning under Temporal Leakage, Distribution Shift, and Market Frictions

\section*{Motivation and Problem}
Financial ML is a useful stress test for reliable data and learning systems. Data are revised, assets enter and leave the observable universe, market regimes change, decisions incur costs, and repeated experimentation makes accidental selection bias likely. Many studies treat these concerns as reporting details. This proposal treats them as coupled systems and methodological problems: how can an experiment ensure that every decision uses only information available at that time, how can a result remain reproducible across data and code revisions, and how can a system prevent weak evidence from being promoted into a decision claim?

\section*{Research Questions}
\begin{enumerate}
\item Which temporal data contracts are necessary to detect leakage introduced by data revisions, universe construction, feature transformation, cross-validation, and model selection?
\item How should evaluation separate predictive quality, decision utility, market frictions, and uncertainty under distribution shift?
\item Which provenance, reconciliation, and monitoring evidence is sufficient for an independent researcher to reproduce or falsify a reported result?
\item Can fail-closed governance gates reduce unsupported claims without making the research workflow unusably rigid?
\end{enumerate}

\section*{Methodology}
\textbf{Work package 1: temporal contract benchmark.} I will define machine-readable contracts for observation time, effective time, revision state, universe membership, and permitted transformations. A benchmark suite will inject known violations into synthetic and versioned real-data pipelines. Detection coverage, false alarms, downstream metric distortion, and decision impact will be measured.

\textbf{Work package 2: shift-aware evaluation.} Models will be trained and selected using rolling or expanding windows with fold-local preprocessing and hyperparameter governance. Tests will include natural regime partitions and controlled perturbations to feature distributions, missingness, noise, and label relationships. Comparisons will include simple statistical and ML baselines. Results will report dispersion across folds and regimes rather than only a pooled average.

\textbf{Work package 3: decision layer with frictions.} Predictive outputs will feed explicit decision rules subject to turnover, transaction-cost, exposure, concentration, and risk constraints. The study will vary cost and liquidity assumptions and will distinguish forecast metrics from decision outcomes. Any economic result will be presented as historical experimental evidence, not as prospective investment advice.

\textbf{Work package 4: reproducible systems and governance.} Each run will preserve code revision, environment, data snapshot identifiers, configuration, random seeds, stage results, and reconciliation checks. Typed stage contracts will prevent downstream use of missing or invalid evidence. A governance layer will issue research-only, observation-only, or hold-for-review outcomes based on predefined checks. Human-subject, legal, and investment-validity claims remain outside scope unless separately established.

\section*{Evaluation}
The evaluation will measure: (i) leak-detection precision and recall on injected failures; (ii) reproducibility of tables and figures from archived evidence; (iii) sensitivity of model and decision rankings to regimes, costs, and reasonable hyperparameter choices; (iv) computational overhead of lineage and checks; and (v) analyst usability through task-based studies if ethics approval and resources permit. Ablations will isolate temporal contracts, monitoring, optimization constraints, and governance gates.

\section*{Expected Contribution and Fit}
The expected contribution is not a universal trading algorithm. It is a set of methods and open infrastructure for more credible research on temporal decision systems. The work would contribute a failure-mode benchmark, explicit data/evaluation contracts, robust decision-evaluation protocols, and an auditable reference architecture.

The proposed fit with Shahin Jabbari is risk-aware optimization and accountable decisions under non-stationarity and competing objectives, grounded in their broader work on machine learning, fairness, optimization, and game theory. At Drexel University, this project could connect financial ML with data systems, optimization, trustworthy AI, or ML systems while retaining a precise empirical core.

\section*{Risks and Mitigations}
Financial data access may limit redistribution; the benchmark will therefore combine public or redistributable sources with synthetic generators and publish transformation code rather than restricted raw data. Market results may be unstable; instability is treated as a finding and evaluated across regimes. System complexity may obscure causal conclusions; preregistered ablations and simple baselines will keep individual claims testable. Any claims about legal compliance, deployment safety, or future financial performance are explicitly excluded.

\section*{Indicative Timeline}
Year 1: literature review, formal problem definitions, temporal-contract benchmark, and baseline pipeline. Year 2: shift-aware learning studies and decision layer. Year 3: reproducible systems architecture, governance experiments, and cross-domain validation. Final period: integrated evaluation, publications, open artifacts, and dissertation.

\section*{Selected Literature to Verify and Expand}
Before submission, replace this section with a supervisor-specific bibliography verified against current papers. Core areas are backtest overfitting, concept drift, time-series validation, data provenance, ML reproducibility, constrained portfolio optimization, and trustworthy ML systems. No unverified citation has been inserted into this draft.
\end{document}
